\section{A subspace intersection inequality}
\label{sec:subspaces}

The algebraic estimate behind the lower bound concerns arbitrary pairs of
subspaces.  The only hypothesis is that each subspace is disjoint from its
designated partner. Polynomial determinants have also been used to
control sums of intersection dimensions in subspace designs
\cite{GuruswamiKopparty2016}. Here the degree filtration is adapted
to the arbitrary family $U_i$, and all details are proved below.

\begin{lemma}[Subspace intersection bound]
\label{lem:subspace}
Let $W$ be a real vector space of dimension $N$, and let
$U_1,\ldots,U_m,V_1,\ldots,V_m$ be subspaces of $W$, where $m\geq1$.
If $U_i\cap V_i=\{0\}$ for every $i$, then
\begin{equation}
\label{eq:subspace-bound}
  \sum_{i,j=1}^{m}\dim(U_i\cap V_j)
  \leq \left\lfloor\frac{m^2}{4}\right\rfloor N.
\end{equation}
\end{lemma}

\begin{proof}
We encode the subspaces by polynomial evaluation constraints.  Evaluating
away from the constraint nodes gives a filtration of evaluation spaces.
At each node, disjointness bounds the evaluated rank.  The resulting
rank deficiency gives a lower bound on a polynomial minor's vanishing
order.  Comparing these orders with the minor's degree, then summing
over the filtration, will give the bound.

Choose distinct real numbers $t_1,\ldots,t_m$.  For a subspace $Z\subseteq W$
and an integer $h\geq0$, define
\[
  \mathcal P_h(Z)=
  \{P\in Z[\xi]:\deg P\leq h,\ P(t_i)\in U_i\text{ for every }i\},
  \qquad \mathcal P_{-1}(Z)=\{0\}.
\]
Here $Z[\xi]$ denotes polynomials whose coefficients belong to $Z$.
Write $E_h(Z;t)=\{P(t):P\in\mathcal P_h(Z)\}$ and put
\[
  k_h(Z)=\dim\mathcal P_h(Z)-\dim\mathcal P_{h-1}(Z),
  \qquad k_{-1}(Z)=0.
\]

\emph{Evaluation away from the nodes.}
For $t\notin\{t_1,\ldots,t_m\}$, the kernel of evaluation at $t$ is
\[
  \ker\bigl(\mathcal P_h(Z)\longrightarrow Z,\ P\longmapsto P(t)\bigr)
  =(\xi-t)\mathcal P_{h-1}(Z).
\]
Indeed, division by $\xi-t$ preserves the coefficient space $Z$, and it
preserves each constraint at $t_i$ because $t_i-t\ne0$.  The converse
inclusion follows by multiplication.  Therefore
\begin{equation}
\label{eq:filtration-evaluation}
  \dim E_h(Z;t)=k_h(Z)
  \qquad(t\notin\{t_1,\ldots,t_m\}).
\end{equation}
Also, evaluation at all $m$ nodes is an isomorphism
\[
  \mathcal P_{m-1}(Z)\longrightarrow
  \bigoplus_{i=1}^{m}(U_i\cap Z).
\]
Its inverse sends $(z_1,\ldots,z_m)$ to the polynomial
$\sum_i L_i(\xi)z_i$, where
$L_i(\xi)=\prod_{k\ne i}(\xi-t_k)/(t_i-t_k)$.
Uniqueness follows because a scalar polynomial of degree at most $m-1$
with $m$ distinct roots vanishes, applied in coordinates of $Z$.
Consequently
\begin{equation}
\label{eq:filtration-sum}
  \sum_{h=0}^{m-1}k_h(Z)
  =\dim\mathcal P_{m-1}(Z)
  =\sum_{i=1}^{m}\dim(U_i\cap Z).
\end{equation}

\emph{Polynomial minors of controlled degree.}
Fix $t_*$ outside the nodes.  The spaces $E_h(W;t_*)$ are nested, so
\[
  \ell_h=k_h(W)-k_{h-1}(W)\geq0
  \qquad(0\leq h\leq m-1).
\]
At stage $h$, choose $\ell_h$ polynomials from $\mathcal P_h(W)$ whose
values at $t_*$ extend the previously chosen values to a basis of
$E_h(W;t_*)$.  In fixed coordinates of $W$, let $M_h(\xi)$ be the matrix
of all polynomial columns chosen through stage $h$.  It has
$r_h=k_h(W)$ columns, independent at $t_*$, and hence has an $r_h$-row
square minor $\Delta_h(\xi)$ with $\Delta_h(t_*)\ne0$.  Each column
introduced at stage $e$ has degree at most $e$.  Expanding the determinant
thus gives
\begin{equation}
\label{eq:minor-degree}
  \deg\Delta_h\leq\sum_{e=0}^{h}e\ell_e.
\end{equation}
If $r_h=0$, use the empty determinant $\Delta_h=1$.

\emph{Rank loss at the nodes.}
For each $j$ we claim that
\begin{equation}
\label{eq:node-rank}
  \dim E_h(W;t_j)\leq k_h(W)-k_h(V_j).
\end{equation}
Let $\pi_j:W\longrightarrow W/V_j$ be the quotient map.  At every
non-node $t$, the kernel of $\pi_j$ on $E_h(W;t)$ contains
$E_h(V_j;t)$.  Equation~\eqref{eq:filtration-evaluation} gives
\[
  \dim\pi_j(E_h(W;t))\leq k_h(W)-k_h(V_j).
\]
To pass to $t_j$, take a fixed basis of the polynomial space
$\mathcal P_h(W)$, evaluate its members at the variable $\xi$, and project
their columns by $\pi_j$.  In fixed coordinates this is a polynomial
matrix.  Every minor larger than the displayed rank bound vanishes at
every non-node, an infinite set, and so is the zero polynomial.  Its rank
at $t_j$ obeys the same bound.  But $E_h(W;t_j)\subseteq U_j$, and $\pi_j$
is injective on $U_j$ because $U_j\cap V_j=\{0\}$.  This proves
\eqref{eq:node-rank}.

The chosen columns of $M_h(t_j)$ lie in $E_h(W;t_j)$, and restricting
to the rows of a minor cannot increase rank.
The square submatrix defining $\Delta_h$ therefore has rank deficiency
at least $k_h(V_j)$ at $t_j$.  A square polynomial matrix with deficiency
at least $q$ at $t_j$ has determinant divisible by $(\xi-t_j)^q$:
choose $q$ independent vectors in its kernel at $t_j$, extend them to
a constant invertible column change, and observe that the corresponding
$q$ transformed columns vanish at $t_j$.  Each contributes the factor
$\xi-t_j$, while the column change rescales the determinant by a nonzero
constant.  Applying this observation separately at the distinct nodes
shows that
\[
  \prod_{j=1}^{m}(\xi-t_j)^{k_h(V_j)}\quad\text{divides }\Delta_h(\xi).
\]
For $r_h=0$, all $k_h(V_j)$ are zero by
\eqref{eq:filtration-evaluation}, so this assertion also holds for the
empty determinant.  Since $\Delta_h$ is not the zero polynomial,
\eqref{eq:minor-degree} now yields
\begin{equation}
\label{eq:multiplicity-budget}
  \sum_{j=1}^{m}k_h(V_j)\leq\sum_{e=0}^{h}e\ell_e.
\end{equation}

\emph{Summing the degree bounds.}
Sum \eqref{eq:multiplicity-budget} over $h=0,\ldots,m-1$ and apply
\eqref{eq:filtration-sum} with $Z=V_j$.  The result is
\begin{align*}
  \sum_{i,j=1}^{m}\dim(U_i\cap V_j)
  &\leq\sum_{h=0}^{m-1}\sum_{e=0}^{h}e\ell_e
   =\sum_{e=0}^{m-1}e(m-e)\ell_e\\
  &\leq\left\lfloor\frac{m^2}{4}\right\rfloor
           \sum_{e=0}^{m-1}\ell_e
   \leq\left\lfloor\frac{m^2}{4}\right\rfloor N.
\end{align*}
The last step uses
$\sum_e\ell_e=k_{m-1}(W)=\dim E_{m-1}(W;t_*)\leq N$.
\end{proof}

\begin{remark}
The constant in Lemma~\ref{lem:subspace} is sharp.  Write $W=P\oplus Q$
and divide the indices into groups of sizes $a=\lfloor m/2\rfloor$ and
$b=\lceil m/2\rceil$.  Set $(U_i,V_i)=(P,Q)$ in the first group and
$(U_i,V_i)=(Q,P)$ in the second.  Only pairs from different groups
contribute, and their total is $ab(\dim P+\dim Q)=\lfloor m^2/4\rfloor N$.
\end{remark}

The form needed for crossing counts follows by taking graphs of linear
maps.  Their intersections identify with kernels of matrix differences, so the
upper bound on intersection dimensions becomes a lower bound on ranks.

\begin{corollary}[Rank sum for differences]
\label{cor:rank-sum}
Let $X_1,\ldots,X_m,Y_1,\ldots,Y_m$ be real $d\times d$ matrices.
If $X_i-Y_i$ is invertible for every $i$, then
\begin{equation}
\label{eq:rank-sum}
  \sum_{\substack{1\leq i,j\leq m\\i\ne j}}
       \rank(X_i-Y_j)\geq2d_m d,
  \qquad
  d_m=\left\lfloor\frac m2\right\rfloor
      \left\lfloor\frac{m-1}{2}\right\rfloor.
\end{equation}
The sum is over ordered pairs of distinct indices.
\end{corollary}

\begin{proof}
In $W=\R^d\oplus\R^d$, let
\[
  U_i=\{(v,X_iv):v\in\R^d\},\qquad
  V_j=\{(v,Y_jv):v\in\R^d\}.
\]
The map $v\mapsto(v,X_iv)$ identifies $\ker(X_i-Y_j)$ with
$U_i\cap V_j$.  Thus
\[
  \dim(U_i\cap V_j)=d-\rank(X_i-Y_j),
  \qquad U_i\cap V_i=\{0\}.
\]
Lemma~\ref{lem:subspace}, with $N=2d$, gives
\[
  \sum_{i,j=1}^{m}\rank(X_i-Y_j)
  \geq\left(m^2-2\left\lfloor\frac{m^2}{4}\right\rfloor\right)d.
\]
The $m$ diagonal terms each equal $d$.  Removing them gives
\eqref{eq:rank-sum}, because
\[
  m^2-m-2\left\lfloor\frac{m^2}{4}\right\rfloor
  =\begin{cases}
      2r(r-1),&m=2r,\\
      2r^2,&m=2r+1,
    \end{cases}
  =2d_m.
\]
\end{proof}
